\documentclass[../../../include/open-logic-section]{subfiles}

\begin{document}

\olfileid{sth}{ordinals}{replacement}
\olsection{الاستبدال}

حفزنا في~\olref[sth][ordinals][vn]{sec} إدخال الأعداد الترتيبية
باقتراح إمكان معاملتها بوصفها أنماط ترتيب، أي ممثلات قياسية للترتيبات
الجيدة. ولكي ينجح ذلك، نحتاج إلى إثبات أن \emph{كل ترتيب جيد متماثل
مع عدد ترتيبي ما}. وهذا يتيح لنا تعريف~$\ordtype{A, <}$ بأنه العدد
الترتيبي~$\alpha$ الذي يحقق
$\tuple{A, <} \isomorphic \alpha$.

لكننا، لسوء الحظ، \emph{لا نستطيع} إثبات النتيجة المنشودة باستعمال
المسلّمات التي أدخلناها إلى الآن فقط. (سنرى السبب في
\olref[replacement][strength]{sec}، أما الآن فالنقطة هي أننا لا نستطيع.)
نحتاج إلى فكرة جديدة، وها هي:

\begin{axiom}[مخطط الاستبدال]
لكل صيغة~$\phi(x, y)$، يكون الآتي مسلّمة:
\begin{quote}
	لكل~$A$، إذا كان~$(\forall x \in A)\lexists![y][\phi(x,y)]$، فإن
	$\Setabs{y}{(\exists x \in A)\phi(x,y)}$ توجد.
\end{quote}
\end{axiom}
\noindent
كما في الفصل، هذا مخطط: فهو يولد عددًا لا نهائيًا من المسلّمات، واحدة
لكل صيغة من الصيغ المختلفة اللانهائية العدد~$\phi$. ويمكن بالقدر نفسه
أن يكتب (وهو المعتاد) هكذا:

\begin{defish}
لكل صيغة~$\phi(x,y)$ لا تحتوي «$B$»، يكون الآتي مسلّمة:
\[
\forall A[(\forall x \in A)\lexists![y][\phi(x,y)] \lif
\exists B\forall y (y \in B \liff (\exists x \in A)\phi(x,y))]
\]
\end{defish}

غير أن هذه، عند اللقاء الأول، صيغة شديدة التشابك. ولعل النتيجة السريعة
الآتية للاستبدال تعبّر بصورة \emph{أوضح} عن الفكرة الحدسية التي نعمل
بها:
\footnote{الحد تعبير يعيّن شيئًا واحدًا بالضبط (عند كل إكمال لمتغيراته
  الحرة). فمثلًا «$\emptyset$» حد يعيّن المجموعة الخالية؛ و«$\{x\}$»
  حد يعيّن المجموعة الأحادية لـ$x$ (أيًا كان~$x$)؛ و«$x \cup y$» حد
  يعيّن اتحاد~$x$ و$y$ (أيا كانا).}

\begin{cor}
لكل حد~$\tau(x)$ ولكل مجموعة~$A$، توجد المجموعة الآتية:
\[
	\Setabs{\tau(x)}{x \in A} = \Setabs{y}{(\exists x \in A)y = \tau(x)}.
\]
\end{cor}

\begin{proof}
لما كانت~$\tau$ \emph{حدًا}، كان
$\forall x \lexists![y][\tau(x) = y]$. وبالأحرى يكون
$(\forall x \in A)\lexists![y][\tau(x) = y]$. ولذلك توجد
$\Setabs{y}{(\exists x \in A)\tau(x) = y}$ بموجب الاستبدال.
\end{proof}
\noindent
يوحي هذا بأن «الاستبدال» اسم حسن للمسلّمة: فإذا أُعطيت مجموعة~$A$،
أمكنك تكوين مجموعة جديدة~$\Setabs{\tau(x)}{x \in A}$، باستبدال كل
عنصر من~$A$ بصورته تحت~$\tau$. بل يمكننا، اقتداء بترميز صورة مجموعة
تحت دالة، أن نكتب~$\funimage{\tau}{A}$ بدلًا من
$\Setabs{\tau(x)}{x \in A}$.

غير أن الأمر الحاسم هو أن~$\tau$ \emph{حد}. ولا يلزم أن يكون (اسمًا
لـ)\emph{دالة} بمعنى~\olref[sfr][fun][rel]{sec}، أي مجموعة معينة من
الأزواج المرتبة. فإذا كانت~$f$ دالة (بذلك المعنى)، فإن المجموعة
$\funimage{f}{A} = \Setabs{f(x)}{x \in A}$ ليست إلا مجموعة جزئية
معينة من~$\ran{f}$، ووجودها مضمون بالفعل باستعمال مسلّمات~$\Zminus$
وحدها.
\footnote{يكفي النظر في
$\Setabs{y \in \bigcup \bigcup f}{(\exists x \in A)y = f(x)}$.}
أما الاستبدال، فهو إضافة \emph{قوية} إلى مسلّماتنا، كما سنرى في
\olref[sth][replacement][]{chap}.


\end{document}
